% ============================================================
% PART VI: CASE STUDY (CHOOSE-ONE SCENARIO)
%
% PEANUT EXAM CASE STUDY LAYOUT STANDARD
%
% AI INSTRUCTION:
% - This part offers several self-contained scenarios; the student picks ONE
% - The instruction box is \fbox{\parbox{0.9\textwidth}{...}} and states:
%     how many to choose, the penalty for choosing more, and the rubric split
% - Each scenario uses \subsubsection*{Scenario A --- <Title>}
% - Every scenario must be answerable in the SAME amount of space and
%   must be of comparable difficulty --- students should not be punished
%   for their choice
% - Structure each scenario as: Context -> sequence of events -> Your Task
% - Give a generous \vspace after the task list (usually 100-140mm)
% - Keep 3 to 5 scenarios
% ============================================================

\section*{Part VI: Case Study \hfill \normalfont\normalsize(N Points)}

\begin{center}
\fbox{%
\parbox{0.9\textwidth}{%
\textbf{INSTRUCTION:} Choose \textbf{exactly ONE} scenario from the options below.
Attempting more than one will result in \textbf{zero points} for this entire part.\\[4pt]
Circle your chosen scenario letter before you begin writing.\\[4pt]
Your answer will be graded on: \textit{correctness of concepts} (N pts),
\textit{justification and reasoning} (N pts), and \textit{clarity of presentation} (N pts).
}}
\end{center}

\vspace{8mm}

% ============================================================
\subsubsection*{Scenario A --- Title of the First Scenario}

\textbf{Context:}
Two or three sentences setting up a concrete system: who runs it,
what it is built from, and what scale it operates at.

The following sequence of events occurs:
\begin{enumerate}
\item The first event (an operator action, a configuration change).
\item The second event (a failure, a spike, an attack).
\item The third event (a recovery, a new resource joining).
\end{enumerate}

\textbf{Your Task:}

\begin{enumerate}[label=(\alph*)]
\item For each event, identify \textbf{which component(s)} detect the change,
make the decision, and take corrective action. Explain the role of each
component in that specific event.

\item Draw a simplified diagram of the system state \textbf{after all events}
(boxes and labels are sufficient).

\item Extend the failure one step further and explain the impact on the
running system. Answer in terms of the architecture taught in this course.
\end{enumerate}

\vspace{120mm}

% ============================================================
\newpage
\subsubsection*{Scenario B --- Title of the Second Scenario}

\textbf{Context:}
A different system with a different failure mode, comparable in difficulty
to Scenario A.

\textbf{Your Task:}

\begin{enumerate}[label=(\alph*)]
\item Analyze the situation using the relevant model from the course.
\item Compute or derive the quantity the scenario depends on.
\item Recommend a remediation and state its trade-off.
\end{enumerate}

\vspace{120mm}

% ============================================================
\newpage
\subsubsection*{Scenario C --- Title of the Third Scenario}

\textbf{Context:}
A third option, typically the more design-flavored or more quantitative one,
so students can play to their strength.

\textbf{Your Task:}

\begin{enumerate}[label=(\alph*)]
\item Identify the core problem and the constraint that makes it hard.
\item Propose a solution and justify each decision.
\item State one thing your solution does \textit{not} handle, and why that is acceptable.
\end{enumerate}

\vspace{120mm}
